\chapter{Exploratory AI-Assisted Measurement}
\label{appendix:ai-assisted-measurement}

This appendix records a small exploratory trial in which an AI coding agent
carried out this study's measurement protocol on its own, over two physics
engines from the selected set (PyBullet and MuJoCo), using the template of
Appendix~\ref{appendix:measurement-template} and the procedure of
Section~\ref{sec:measurement-procedure}. It is not part of the validated results
in Chapter~\ref{ch:measurement-results}; it is a feasibility check and the basis
for the future-work direction in Section~\ref{sec:future-work}. The agent worked
in a Linux shell with file-system and network access in early June 2026, so its
repository snapshots fall slightly later than the May 2026 snapshot of the main
study.

\section*{Task Given to the Agent}

The agent was given a single natural-language prompt and no installation commands;
it had to locate the official documentation, install each package, verify the
installation, and complete the template unaided. The prompt is reproduced below.

\begin{quote}
\small
Read the CSV file in the current directory and use the first column
(``Metric'') as the measurement template. Measure two software packages,
PyBullet and MuJoCo. The requirements were:
\begin{itemize}[noitemsep]
  \item Only fill in the columns for these two packages; do not modify the other
        software columns. A new CSV file may be created.
  \item Do not rely on user-provided installation commands; find the official
        documentation, GitHub README, or project website to determine the
        installation and verification steps.
  \item Follow each metric in the CSV to complete installation, verification,
        observation, and data entry.
  \item Take a screenshot at every step and save it to a screenshots folder.
        Screenshots must cover at least: reading the CSV, finding the official
        documentation, the installation steps, the verification steps, error
        messages or successful output, and the final CSV result.
  \item Evidence may come only from official documentation, GitHub, terminal
        output, and commands actually executed.
  \item Write ``unclear'' for uncertain values and ``n/a'' for values that do
        not apply.
  \item Also write a \texttt{notes.md} file recording the selected packages, the
        official links used, the commands run at each step, the screenshot
        filenames, whether installation succeeded, whether verification
        succeeded, which CSV cells were filled or changed, and any problems
        encountered with their recovery.
\end{itemize}
Complete Bullet/PyBullet first, then MuJoCo.
\end{quote}

\section*{Official Sources Located by the Agent}

The agent found the official sources in Table~\ref{tab:ai-pilot-sources} and used
them to choose the installation and verification steps.

\begin{table}[H]
\centering
\caption{Official sources the agent used for the two packages.}
\label{tab:ai-pilot-sources}
\small
\begin{tabular}{@{}p{2.6cm}p{3.8cm}p{6.2cm}@{}}
\toprule
\textbf{Package} & \textbf{Resource} & \textbf{Location} \\
\midrule
PyBullet & Website and forum  & \url{https://pybullet.org} \\
         & Source repository  & \url{https://github.com/bulletphysics/bullet3} \\
         & Python package     & \url{https://pypi.org/project/pybullet/} \\
\midrule
MuJoCo   & Website            & \url{https://mujoco.org} \\
         & Source repository  & \url{https://github.com/google-deepmind/mujoco} \\
         & Documentation      & \url{https://mujoco.readthedocs.io} \\
\bottomrule
\end{tabular}
\end{table}

\section*{Installation and Verification}

Both packages installed with a single \texttt{pip} command inside a Python
virtual environment, and both were verified with a short program that the agent
wrote and ran. Table~\ref{tab:ai-pilot-outcome} summarizes the outcome.

\begin{table}[H]
\centering
\caption{Installation and verification outcome of the AI-assisted trial.}
\label{tab:ai-pilot-outcome}
\small
\begin{tabular}{@{}p{3.1cm}p{4.7cm}p{4.7cm}@{}}
\toprule
\textbf{Item} & \textbf{PyBullet} & \textbf{MuJoCo} \\
\midrule
Install command     & \texttt{pip install pybullet} & \texttt{pip install mujoco} \\
Installed version   & 3.2.7 & 3.9.0 \\
Installation steps  & 1 (single pip command) & 1 (single pip command) \\
Extra packages      & none (self-contained) & bundled with the wheel \\
Operating system    & Linux (Ubuntu) & Linux (Ubuntu) \\
Verification        & Loaded the bundled \texttt{plane.urdf} model in DIRECT mode & Compiled an inline XML model of a free-jointed box on a plane \\
Result              & Succeeded & Succeeded \\
\bottomrule
\end{tabular}
\end{table}

The agent also queried the GitHub API and recorded the template's project
metadata, each entry tied to the source it read. For PyBullet
(\path{bulletphysics/bullet3}) it recorded 311 contributors, a zlib license, a
first commit in April 2011, and a most recent push in October 2025; for MuJoCo
(\path{google-deepmind/mujoco}), 116 contributors, an Apache-2.0 license, a
first commit in August 2021, and a most recent push in June 2026. These
contributor counts come from the GitHub contributors API, which returns fewer
entries than the commit-log method of the main study
(Section~\ref{subsec:repository-metrics}; 333 developers for Bullet/PyBullet,
138 for MuJoCo), so the two are not directly comparable.

\section*{Problems Encountered and Recovery}

The agent encountered three problems and recovered from each without manual
intervention. First, a direct \texttt{pip install} failed because the system
Python was externally managed; the agent created a virtual environment and
installed into it. Second, its first PyBullet verification script queried the
physics engine before connecting to it; the agent reordered the calls so the
connection came first. Third, no screenshot utility was present, so the agent
installed one.

\section*{Limitations and Relationship to the Main Study}

As a single run over two packages by one agent, this shows feasibility, not a
result. The checkable, factual entries (installation command, installed version,
license, contributor count, verification result) were each tied to a public
source and straightforward to confirm. The subjective impression scores and
free-text judgements were not validated here, and in places disagree with the
agent's own notes, which is the clearest reason a human must still review the
output before it can serve as evidence.

The trial's artifacts, the completed template and the notes document, are
preserved in the supplementary project materials. The trial motivates the
agent-assisted measurement under human supervision discussed in
Section~\ref{sec:future-work}.
